%% ============================================================
%%  CV - Kossi Samuel Gabiam (Version Française Haute Qualité)
%%  Compile: pdflatex cv_gabiam_fr.tex
%% ============================================================
\documentclass[11pt, a4paper]{article}

% ── Packages ──────────────────────────────────────────────
\usepackage[T1]{fontenc}
\usepackage[utf8]{inputenc}
\usepackage[french]{babel}
\usepackage{mathptmx} % Times New Roman
\usepackage[margin=1.5cm, top=1.5cm, bottom=1.5cm]{geometry}
\usepackage{hyperref}
\usepackage{enumitem}
\usepackage{titlesec}
\usepackage{tabularx}
\usepackage{array}
\usepackage{xcolor}
\usepackage{parskip}
\usepackage{fancyhdr}
\usepackage{needspace}

% ── Couleurs ──────────────────────────────────────────────
\definecolor{sectioncolor}{HTML}{1E40AF} % Bleu Marine Élégant (Grands Titres)
\definecolor{secondary}{HTML}{475569}    % Gris Sombre (Numérotation & Métadonnées)
\definecolor{linkcolor}{HTML}{1E40AF}    % Bleu pour les Liens

% ── Hyperliens ────────────────────────────────────────────
\hypersetup{
  colorlinks=true,
  urlcolor=linkcolor,
  linkcolor=linkcolor
}

% ── Numérotation de page ──────────────────────────────────
\pagestyle{fancy}
\fancyhf{}
\renewcommand{\headrulewidth}{0pt}
\fancyfoot[C]{\small\color{secondary}\thepage}

% ── Style des Sections ────────────────────────────────────
\titleformat{\section}
  {\large\bfseries\color{sectioncolor}\scshape}
  {}{0em}
  {}
  [\vspace{-4pt}\color{sectioncolor}\rule{\linewidth}{1pt}]
\titlespacing{\section}{0pt}{14pt}{6pt}

% ── Configuration des Listes ──────────────────────────────
\setlist[itemize]{
  leftmargin=1.3em,
  itemsep=3pt,
  topsep=2pt,
  parsep=0pt,
  label=\small\textbullet
}

% ── Macro Entrée CV ───────────────────────────────────────
\newcommand{\cventry}[4]{%
  \needspace{3\baselineskip}%
  \noindent\textbf{\large #1} \hfill {\small\color{secondary}\textit{#2}}\\[2pt]%
  \ifx&#3&\else\noindent\textbf{\textit{#3}}\\[3pt]\fi
  \ifx&#4&\else\noindent{#4}\fi
  \vspace{6pt}\par
}

% =============================================================
\begin{document}
\color{black} % Texte noir strictement dominant dans tout le document

% ── EN-TÊTE ───────────────────────────────────────────────────
\begin{center}
  {\Huge\bfseries\scshape Kossi Samuel Gabiam}\\[6pt]
  {\large\bfseries Développeur Full-Stack \textbullet{} Ingénieur Systèmes \& Cloud}\\[8pt]
  \small
  \href{mailto:gabiam.k.samuel@gmail.com}{gabiam.k.samuel@gmail.com}
  \quad\textbullet\quad
  +228 90 84 77 68
  \quad\textbullet\quad
  Lomé, Togo\\[3pt]
  \href{https://www.linkedin.com/in/samuel-g-107663165/}{linkedin.com/in/samuel-g-107663165}
  \quad\textbullet\quad
  \href{https://github.com/Shieldish}{github.com/Shieldish}
  \quad\textbullet\quad
  \href{https://portfolio.gabiam-k-samuel.workers.dev/}{portfolio.gabiam-k-samuel.workers.dev}\\[4pt]
  {\small\color{secondary}\textit{Disponible Partout : Télétravail, Présentiel \& Déplacement}}
\end{center}

\vspace{-2pt}
\noindent\rule{\linewidth}{1pt}
\vspace{6pt}

% ── PROFIL ────────────────────────────────────────────────────
\section{Profil Professionnel}

Développeur Full-Stack expérimenté, titulaire d'un Master Professionnel en Systèmes, Réseaux \& Cloud Computing et d'une Licence en Génie Logiciel de la Faculté des Sciences de Sfax (Tunisie). Spécialisé dans la conception d'architectures distribuées hautement disponibles, le développement d'applications web/mobiles complexes et la mise en place d'infrastructures Cloud et pipelines DevOps.

% ── EXPÉRIENCE PROFESSIONNELLE ────────────────────────────────
\section{Expérience Professionnelle}

\cventry{Développeur Full-Stack}{Juin 2025 - Présent}{SuiSco Sarl}{Lomé, Togo}
\begin{itemize}
  \item \textbf{EduBoost (\href{https://www.eduboost.tg}{eduboost.tg}) :} Architecture et développement complet de la plateforme nationale de tombola multi-canaux (Web, SMS, USSD, Mobile Money MIXX/VAS) actuellement en production au Togo.
  \item \textbf{Architecture Distribuée :} Conception d'un système hautement disponible à 15 nœuds Django d'écriture avec génération de tickets anti-collision (\texttt{SERVER\_ID} + séquence PostgreSQL) et 1 nœud central d'agrégation.
  \item \textbf{Flux Asynchrones \& Performance :} Synchronisation Celery $\sim$500 ms des tickets vers la base centrale, tâche de crash-recovery (30 s), génération Web asynchrone (polling \texttt{request\_id}) et SMS/USSD synchrone ($<$50 ms).
  \item \textbf{Portail Offrir \& Sécurité :} Déploiement du portail entreprise (génération PDF, OTP SMS/email via route centrale sécurisée, claim), vérification HMAC-SHA256 des webhooks et rate-limiting Redis.
  \item \textbf{Système de Tirage Live :} Définition du workflow en 4 phases avec anonymisation bijective des tickets, vote OTP (\href{https://soiree.eduboost.tg}{soiree.eduboost.tg}), sélection admin et écran de tirage public (\href{https://tirage.eduboost.tg}{tirage.eduboost.tg}).
  \item \textbf{Projets Complémentaires :} Développement de la plateforme \textbf{DefiQuiz} (\href{https://defiquiz.tg/}{defiquiz.tg}), maintenance de l'application \textbf{SuiSco My Account} (PHP Laravel) et intégration du Single Sign-On centralisé \textbf{Keycloak} (OAuth2/OIDC).
\end{itemize}

% ── PROJETS MAJEURS ───────────────────────────────────────────
\needspace{6\baselineskip}
\section{Projets Majeurs}

\cventry{EduBoost - Plateforme Nationale de Tombola}{2025 - Présent}{SuiSco Sarl / YAS TOGO}{Plateforme en Production}
\begin{itemize}
  \item \textbf{Architecture \& Échelle :} 15 nœuds d'écriture Django + 1 nœud central d'agrégation (7 applications, 3 domaines publics).
  \item \textbf{Multi-Canaux d'Achat :} Web (asynchrone), SMS/USSD (synchrone $<$50 ms, idempotent par \texttt{ref\_id}) et Mobile Money (MIXX \& VAS avec webhooks HMAC-SHA256).
  \item \textbf{Portail Entreprise Offrir :} Gestion des cadeaux d'entreprise, génération de reçus PDF et validation d'OTP centralisée.
  \item \textbf{Système de Tirage Live :} Anonymisation bijective pré-tirage, portail de vote OTP et écran d'affichage du tirage en direct.
  \item \textbf{Technologies Utilisées :} Django 6, DRF, SimpleJWT, Nuxt 3, Vue 3, Pinia, PostgreSQL 17, Redis 7, Celery, Gunicorn, Docker, Tailwind.
  \item \textbf{Liens :} \href{https://www.eduboost.tg}{eduboost.tg} \quad\textbullet\quad \href{https://soiree.eduboost.tg}{soiree.eduboost.tg} \quad\textbullet\quad \href{https://tirage.eduboost.tg}{tirage.eduboost.tg}
\end{itemize}

\vspace{4pt}

\cventry{BudaxDrive - Plateforme d'Autopartage \& Car-Sharing}{2025 - 2026}{Projet Full-Stack \& Cloud}{Plateforme en Production}
\begin{itemize}
  \item \textbf{Écosystème Multi-Plateformes :} Architecture et développement complet d'une solution d'autopartage (Burundi) comprenant une API Spring Boot 3, un portail web Next.js 16 et une application mobile Flutter.
  \item \textbf{Architecture Backend \& Données :} API Java 17 organisée en \textit{Vertical Slice DDD}, 20 migrations Flyway (PostgreSQL), indexation spatiale, cache et rate-limiting distribué avec Redis.
  \item \textbf{Temps Réel Bicanal (WebSocket) :} Broker STOMP bicanal (\texttt{/ws} avec SockJS pour le web, \texttt{/ws-native} TCP pour Flutter) avec authentification JWT au handshake pour la messagerie et les alertes.
  \item \textbf{Sécurité \& Audit RGPD :} Refonte du cycle de vie des documents légaux (carte grise, assurance, contrôle technique) avec streaming sécurisé d'octets d'images, logs d'audit d'accès et soft-delete pour le droit à l'oubli.
  \item \textbf{Inspection Numérique Anti-Fraude :} Protocole d'état des lieux Check-in/Check-out géolocalisé (photos obligatoires, validation PIN et coordonnées GPS via l'application mobile).
  \item \textbf{Technologies Utilisées :} Java 17, Spring Boot 3.2, PostgreSQL, Redis, Flyway, Next.js 16, React 19, TypeScript, Flutter, Dart, Riverpod, Docker, Fly.io, Vercel.
  \item \textbf{Liens \& Déploiement :} \href{https://carshare-frontend.vercel.app/}{carshare-frontend.vercel.app} (Web Live) \quad\textbullet\quad GitLab : \href{https://gitlab.com/orlyndonse/carshare-backend}{Backend} / \href{https://gitlab.com/orlyndonse/carshare-frontend}{Frontend} / \href{https://gitlab.com/orlyndonse/carshare-mobile-flutter}{Mobile}
\end{itemize}

\vspace{4pt}

\cventry{GTEF Platform — Monorepo EdTech \& E-Learning (LMS)}{2025 - 2026}{Global Tech Education Foundation}{Monorepo Turborepo}
\begin{itemize}
  \item \textbf{Architecture Monorepo :} Conception d'une infrastructure unifiée avec Turborepo et pnpm comprenant 3 applications (Site institutionnel/LMS, Portail d'administration ERP/CRM, API haute performance Fastify 5) et 4 packages partagés (\texttt{database}, \texttt{ui}, \texttt{i18n}, \texttt{config}).
  \item \textbf{API Backend \& Recherche :} Serveur Fastify 5 modulaire avec Drizzle ORM sur PostgreSQL (Neon), recherche plein texte temps réel avec Meilisearch et documentation OpenAPI dynamique via Scalar.
  \item \textbf{E-Learning \& Moteur de Certification :} Suivi de progression pédagogique, quiz interactifs, génération dynamique de certificats PDF horodatés avec QR Code cryptographique (\texttt{pdf-lib}, \texttt{qrcode}).
  \item \textbf{Multi-Paiements \& Cloud :} Double passerelle de paiement intégrée (Stripe et PayDunya pour le Mobile Money africain), stockage d'actifs pédagogiques sur Cloudflare R2 / S3 et observabilité Sentry.
  \item \textbf{Design System \& Qualité :} Interface moderne Next.js 15 (React 19, Tailwind 4, Motion) bilingue FR/EN, composants réutilisables partagés, tests E2E Playwright et tests unitaires Vitest avec pipeline CI/CD GitHub Actions.
  \item \textbf{Technologies Utilisées :} TypeScript, Fastify 5, Drizzle ORM, PostgreSQL, Neon, Meilisearch, Next.js 15, React 19, Tailwind CSS 4, Turborepo, Docker, Playwright, Vitest.
  \item \textbf{Liens \& Déploiement :} \href{https://gtef-platform-web.vercel.app/}{gtef-platform-web.vercel.app} (Web Live) \quad\textbullet\quad \href{https://github.com/Shieldish/gtef-platform}{Code Source (GitHub)}
\end{itemize}

\vspace{4pt}

\cventry{DefiQuiz - Plateforme de Quiz Éducatif}{2025 - Présent}{SuiSco Sarl / YAS TOGO}{Plateforme Web}
\begin{itemize}
  \item \textbf{Description :} Application de quiz multi-clients en ligne proposant plus de 5 000 questions dans divers domaines de connaissances.
  \item \textbf{Stack Technique :} Frontend Nuxt.js / Next.js, Backend Hono.js / Deno.js.
  \item \textbf{Lien :} \href{https://defiquiz.tg/}{defiquiz.tg}
\end{itemize}

\vspace{4pt}

\cventry{SuiSco My Account \& SSO Centralisé}{2025}{SuiSco Sarl}{Applications Internes}
\begin{itemize}
  \item \textbf{SuiSco My Account :} Application web multi-clients de gestion de comptes développée en PHP Laravel avec API REST et MySQL.
  \item \textbf{SSO Keycloak :} Mise en place d'une solution d'authentification unique (Single Sign-On) basée sur Keycloak (OAuth2/OIDC) pour sécuriser l'accès aux microservices.
\end{itemize}

\vspace{4pt}

\cventry{Plateforme de Gestion de Stages}{2024}{Djagora}{Projet de Fin d'Études}
\begin{itemize}
  \item \textbf{Description :} Application web et mobile pour la gestion des offres de stage, des candidatures et du suivi des étudiants.
  \item \textbf{Technologies :} React, Node.js, MySQL (Web) et React Native (Mobile).
  \item \textbf{Code Source :} \href{https://github.com/Shieldish/PFE-}{github.com/Shieldish/PFE-} \quad\textbullet\quad \href{https://github.com/Shieldish/Gestion-Apps-final}{github.com/Shieldish/Gestion-Apps-final}
\end{itemize}

\vspace{4pt}

\cventry{Extraction \& Visualisation de Données Sociales}{2021}{DK-Soft}{Projet de Stage}
\begin{itemize}
  \item \textbf{Description :} Application web pour l'extraction et la visualisation analytique en temps réel de données des réseaux sociaux.
  \item \textbf{Technologies :} Angular, Elasticsearch, Logstash, Kibana (Stack ELK).
  \item \textbf{Code Source :} \href{https://github.com/Shieldish/test}{github.com/Shieldish/test} \quad\textbullet\quad \href{https://github.com/Shieldish/Api-test}{github.com/Shieldish/Api-test}
\end{itemize}

% ── FORMATION ────────────────────────────────────────────────
\needspace{6\baselineskip}
\section{Formation Académique}

\cventry{Master Professionnel en Systèmes, Réseaux \& Cloud Computing}{2022 - 2024}{Faculté des Sciences de Sfax}{Sfax, Tunisie}
\begin{itemize}
  \item Ingénierie des Systèmes Informatiques - Architectures distribuées, virtualisation, cloud computing, administration système Linux et sécurité réseau.
\end{itemize}

\vspace{4pt}

\cventry{Licence Nationale en Science de l'Informatique}{2019 - 2022}{Faculté des Sciences de Sfax}{Sfax, Tunisie}
\begin{itemize}
  \item Spécialité Génie Logiciel \& Systèmes d'Information - Algorithmique, bases de données, UML, génie logiciel et réseaux informatiques.
\end{itemize}

\vspace{4pt}

\cventry{Baccalauréat Scientifique - Série D}{2014 - 2017}{Lycée La Grâce}{Sokodé, Togo}{}

% ── COMPÉTENCES TECHNIQUES ────────────────────────────────────
\needspace{6\baselineskip}
\section{Compétences Techniques}

\begin{tabularx}{\linewidth}{@{} >{\bfseries}l @{\hspace{14pt}} X @{}}
  Frontend      & React, Next.js, Vue.js, Nuxt 3, Angular, TypeScript, Tailwind CSS \\[4pt]
  Backend       & Django, DRF, Celery, Gunicorn, Node.js, PHP Laravel, Hono.js, Deno.js, API REST, Python \\[4pt]
  Mobile        & React Native, Flutter \\[4pt]
  Bases de données & PostgreSQL, MySQL, Redis, MongoDB, Elasticsearch \\[4pt]
  Systèmes \& Cloud & Linux/Unix, Docker, Docker Compose, Architectures distribuées, Keycloak (SSO), ELK \\[4pt]
  DevOps \& Méthodes & Git, GitHub, GitLab, Pipelines CI/CD, Virtualisation, Agile / Scrum \\
\end{tabularx}

% ── LANGUES ──────────────────────────────────────────────────
\vspace{6pt}
\section{Langues}

\begin{tabularx}{\linewidth}{@{} >{\bfseries}l @{\hspace{14pt}} X @{}}
  Français  & Langue maternelle \\[3pt]
  Anglais   & Compétence professionnelle \\
\end{tabularx}

\end{document}

